\documentclass{beamer}
% Overlay specifications. Each frame's expected page count is written in
% its comment; the deck must produce 3 + 3 + 2 + 2 + 3 + 1 = 14 pages.
\setbeamertemplate{navigation symbols}{}

\title{Overlays}
\author{R. Okafor}
\date{March 2026}

\begin{document}

% 3 slides: two \pause commands inside the list.
\begin{frame}
  \frametitle{Pause}
  Three steps, revealed one at a time.
  \begin{itemize}
    \item Parse the source into tokens.
    \pause
    \item Expand macros and build paragraphs.
    \pause
    \item Break lines and pages.
  \end{itemize}
\end{frame}

% 3 slides: the largest overlay number mentioned is 3.
\begin{frame}
  \frametitle{Item overlays}
  \begin{itemize}
    \item<1-> Always visible from the first slide.
    \item<2-> Appears on the second slide and stays.
    \item<3-> Appears on the third slide.
    \item<2-> Also from the second slide onwards.
  \end{itemize}
\end{frame}

% 2 slides: \only removes the other branch entirely, so the paragraph
% below moves up on slide 1 relative to slide 2.
\begin{frame}
  \frametitle{Only}
  \only<1>{On the first slide this sentence is typeset and the alternative
  text is absent from the page.}%
  \only<2>{On the second slide the alternative text is typeset instead, and
  it is longer so that it takes one more line than the first version did.}

  This closing paragraph is on both slides.
\end{frame}

% 2 slides: \uncover keeps the space on slide 1; \alert colours on slide 2.
\begin{frame}
  \frametitle{Uncover and alert}
  The engine is \alert<2>{incremental}: it keeps the layout of every
  unchanged paragraph.

  \uncover<2->{This paragraph is invisible on the first slide but its space
  is reserved, so the sentence below does not move between the two slides.}

  The line after the uncovered paragraph.
\end{frame}

% 3 slides: \onslide<2-> and \onslide<3-> in sequence.
\begin{frame}
  \frametitle{Onslide}
  First part, on every slide.
  \onslide<2->
  Second part, from the second slide.
  \begin{itemize}
    \item A bullet that comes with the second part.
  \end{itemize}
  \onslide<3->
  Third part, from the third slide.
\end{frame}

% 1 slide: no overlay specification at all.
\begin{frame}
  \frametitle{No overlays}
  A frame without overlay specifications produces exactly one page.
\end{frame}

\end{document}
